\chapter{Progress Yang Telah Dicapai}

\section{Modul Autentikasi dan Multi-tenant}

Progress pada modul autentikasi dan multi-tenant meliputi:

\begin{itemize}
    \item Implementasi autentikasi tiga skema, yaitu API Key untuk komunikasi klien ke gateway, JWT Bearer untuk dashboard, dan Verify Token untuk komunikasi gateway ke agent klien.
    \item Pendaftaran klien atau tenant baru beserta pembuatan API Key otomatis.
    \item Pemisahan akses agar setiap tenant dapat memiliki konfigurasi dan data audit masing-masing.
\end{itemize}

\section{Modul Ingestion}

Progress pada modul ingestion meliputi:

\begin{itemize}
    \item Ingestion data melalui Kafka CDC per klien dengan reconciliation otomatis setiap 15 detik.
    \item Consumer Kafka dapat menerima event perubahan database dan meneruskannya ke proses normalisasi audit log.
\end{itemize}

\section{Modul Kriptografi dan Blockchain}

Progress pada modul kriptografi dan blockchain meliputi:

\begin{itemize}
    \item Hashing SHA3-256 terhadap delapan field kanonik log audit.
    \item Pembangunan Merkle Tree dengan dukungan proof chain lengkap dari leaf ke root.
    \item Dukungan kasus leaf ganjil melalui mekanisme self-pairing.
    \item Anchoring Merkle Root ke jaringan Hyperledger Fabric.
\end{itemize}

\section{Modul Verifikasi}

Progress pada modul verifikasi meliputi:

\begin{itemize}
    \item Verifikasi empat lapis, yaitu eksistensi database, re-hashing lokal, verifikasi ke agent klien, dan pencocokan Merkle Proof terhadap ledger Fabric.
    \item Dashboard menampilkan status integritas granular, seperti valid, tampered, pending, dan unreachable.
    \item Auditor dapat melihat hasil verification untuk mengetahui apakah evidence masih sesuai dengan hash dan proof yang tersimpan.
\end{itemize}

\section{Dashboard dan Admin Panel}

Progress pada dashboard dan admin panel meliputi:

\begin{itemize}
    \item Statistik ringkasan, seperti total log, pending, dan anchored.
    \item Riwayat transaksi dengan pagination, filter, dan pencarian.
    \item Manajemen klien, konfigurasi Kafka, dan konfigurasi agent lapis 3.
    \item Panel monitoring CDC user klien.
\end{itemize}

\placeholderimage{Screenshot dashboard utama Audit Logs atau Overview}{Masukkan screenshot dashboard utama yang menunjukkan ringkasan log, status, dan tabel audit.}

\placeholderimage{Screenshot Admin Panel Client Registry}{Masukkan screenshot halaman admin untuk registrasi client dan pengelolaan konfigurasi tenant.}

\section{Ringkasan Status Modul}

\begin{longtable}{p{4.2cm}p{2.7cm}p{6.8cm}}
\toprule
\textbf{Modul} & \textbf{Status} & \textbf{Keterangan} \\
\midrule
Autentikasi dan Multi-tenant & Berjalan & API Key, JWT, Verify Token, dan registrasi tenant sudah tersedia. \\
\midrule
Kafka CDC Ingestion & Berjalan & CDC per client sudah masuk melalui Debezium dan Kafka. \\
\midrule
Hashing dan Merkle Tree & Berjalan & SHA3-256 dan Merkle proof sudah menjadi bagian core evidence. \\
\midrule
Fabric Anchoring & Berjalan & Merkle Root dapat di-anchor ke Hyperledger Fabric. \\
\midrule
Dashboard Monitoring & Berjalan & Audit logs, client registry, verification, dan monitoring user tersedia. \\
\midrule
Actor Tracking & Dalam penyempurnaan & Mapping actor per operasi sudah masuk ke form, integrasi pipeline masih dilanjutkan. \\
\bottomrule
\end{longtable}
